\section{The unchanged canonical ratio section in the actual native action lift}

\begin{theorem}[The actual native action lift preserves the original local ratio section]
\label{thm:pl-original-native-action-lift-canonical-ratio-section}
\lean{EnvelopeProperty.ProjSubvariety.normalizedSimultaneousNativeCartierFactors_isInvertible,EnvelopeProperty.ProjSubvariety.normalizedSimultaneousNoncentralOpen_off_nativeCartierFactors,EnvelopeProperty.ProjSubvariety.normalizedSimultaneousNativeActionLift_ratioLocalSection}
\leanok
\uses{thm:pl-original-normalized-native-test-configuration,thm:pl-original-normalized-cartier-factor-linearizations,cor:original-pl-group-canonical-ratio-section}
For the actual constructed original normalized test configuration, the
retained canonical ratio group comparison, pulled back along its actual
noncentral action lift, carries the canonical local ratio section along
the action composite to the same canonical local ratio section along
the projection composite. The first composite is the unchanged original
noncentral-open inclusion; the second factors through the actual original
fibre at $1$ and its unchanged noncentral-open factorization.
\end{theorem}
\begin{proof}
The two actual normalized original ideal factors are Cartier. Their
supports avoid the whole original noncentral open, because the original
ideals contain parameter powers. The action composite of the actual
native lift equals its original open inclusion. Its projection composite
equals the original fibre-one projection followed by the original fibre
inclusion, and this latter inclusion factors through the same original
noncentral open. Thus both actual composites factor through the same
open off the numerator support. The canonical Cartier ratio transport
theorem applies to these exact two factorizations and the retained
inclusion-preserving original factor comparisons. Its inverse-numerator
and evaluation-dual sections are exactly the existing original local
Cartier ratio section. This gives the entire claimed module-morphism
equality before taking a residue fibre, valuation or model norm.
\end{proof}
